%% kapa: exam=2 round=30 year=2019 period=1 subject=appraisal_practice subject_ko=감정평가실무 seq=1 number=1 points=40
%% source: Q-Net 감정평가사 기출문제 (cst00309) artlSeq=5205595 file=raw/1008/5205595/감정평가실무.pdf
%% conversion: kordoc 4.24.51 -> tools/gen2.py (auto); review=proofread-sonnet; flags=none
\begin{longquestion}{【문제 1】}
감정평가사 甲은 식료품 제조업을 영위하는 (주)A로부터 일반거래(시가참고) 목적의 감정평가를 의뢰받았다. 관련법규 및 이론을 참작하고 제시된 자료를 활용하여 다음의 물음에 답하시오.(단, 기준시점은 2020. 01. 01임) (40점) \par
\medskip\noindent 물음 1) (주)A의 기업가치를 평가하시오. (25점) \par
\medskip\noindent 물음 2) (주)A의 특허권의 유효 잔존수명을 산출하고, 특허권 가치를 평가하시오. (10점) \par
\medskip\noindent 물음 3) (주)A의 영업권 가치를 평가하시오. (5점) \par
< 자료 1 > 대상 기업 및 특허권 개요 \par
1. 대상 기업 현황 \par
{\small\centering
\begin{tabular}{|c|c|}
\hline
\makecell{상호} & \makecell{(주)A} \\
\hline
\makecell{대표자} & \makecell{이ㅇㅇ} \\
\hline
\makecell{설립일자} & \makecell{2012. 06. 17} \\
\hline
\makecell{사업자번호} & \makecell{514-87-*****} \\
\hline
\makecell{주요제품} & \makecell{과자류} \\
\hline
\end{tabular}\par}

※ 대상 기업은 식료품 제조업을 영위함 \par
2. 특허권 개요 \par
{\small\centering
\begin{tabular}{|c|c|}
\hline
\makecell{명칭} & \makecell{나선형 ** 코팅 장치} \\
\hline
\makecell{등록번호} & \makecell{10-13*****} \\
\hline
\makecell{출원일} & \makecell{2013. 05. 26} \\
\hline
\makecell{특허권자} & \makecell{(주)A} \\
\hline
\makecell{존속기간 만료일} & \makecell{2033. 05. 26} \\
\hline
\end{tabular}\par}

< 자료 2 > 주요가정 \par
1. 추정기간이란 할인현금수지분석법 적용에 있어 현금흐름을 직접 추정하는 기간으로 대상 기업의 특성 및 시장상황 등을 고려하여 5년(1기∼5기)으로 가정함 \par
2. 추정기간이 지난 6기부터는 성장율 0 \%를 적용하며, 5기의 현금흐름이 지속되는 것으로 가정함 \par
3. 대상 기업의 결산일은 매년 말일이며, 현금흐름은 편의상 기말에 발생하는 것으로 가정함 \par
4. 대상 기업의 현금흐름 추정시 비영업용 자산에 의한 수익, 비용은 제외된 것으로 가정함 \par
< 자료 3 > 재무상태표 및 손익계산서 일부 발췌 (2019. 12. 31. 현재) \par
1. 재무상태표(일부 발췌) \par
{\small\centering
\begin{tabular}{|c|c|}
\hline
\makecell{계정과목} & \makecell{금액(원)} \\
\hline
 &  \\
\hline
\makecell{자산} &  \\
\hline
\makecell{Ⅰ.유동자산} &  \\
\hline
\makecell{1.당좌자산} &  \\
\hline
\makecell{(1)단기금융상품} & \makecell{700,000,000} \\
\hline
\makecell{(2)그 외} & \makecell{500,000,000} \\
\hline
\makecell{2.재고자산 등} & \makecell{600,000,000} \\
\hline
 &  \\
\hline
\makecell{Ⅱ.비유동자산} &  \\
\hline
\makecell{1.투자자산} &  \\
\hline
\makecell{(1)장기투자자산} & \makecell{300,000,000} \\
\hline
\makecell{2.유형자산} &  \\
\hline
\makecell{(1)토지} & \makecell{2,500,000,000} \\
\hline
\makecell{(2)건물} & \makecell{1,000,000,000} \\
\hline
\makecell{(3)기계} & \makecell{800,000,000} \\
\hline
 &  \\
\hline
\makecell{부채} &  \\
\hline
\makecell{Ⅰ.유동부채} & \makecell{1,100,000,000} \\
\hline
\makecell{Ⅱ.비유동부채} &  \\
\hline
\makecell{1.장기차입금} & \makecell{2,500,000,000} \\
\hline
\end{tabular}\par}

※ 대상 기업의 무형자산은 영업권과 특허권만 존재함 \par
※ 대상 기업의 비영업용 항목은 단기금융상품, 장기투자자산임 \par
2. 손익계산서(일부 발췌) \par
(단위: 원) \par
{\small\centering
\begin{tabular}{|c|c|c|c|}
\hline
\makecell{구분} & \makecell{2017년} & \makecell{2018년} & \makecell{2019년} \\
\hline
\makecell{매출액} & \makecell{2,000,000,000} & \makecell{2,100,000,000} & \makecell{2,205,000,000} \\
\hline
\makecell{매출원가} & \makecell{1,000,000,000} & \makecell{1,050,000,000} & \makecell{1,102,500,000} \\
\hline
\makecell{매출총이익} & \makecell{1,000,000,000} & \makecell{1,050,000,000} & \makecell{1,102,500,000} \\
\hline
\makecell{판매비와 관리비} & \makecell{200,000,000} & \makecell{210,000,000} & \makecell{220,500,000} \\
\hline
\makecell{영업이익} & \makecell{800,000,000} & \makecell{840,000,000} & \makecell{882,000,000} \\
\hline
\end{tabular}\par}

< 자료 4 > 재무제표 관련 추가 자료 \par
1. 추정기간 동안 매출액은 다음에서 산출한 증가율과 동일한 비율로 증가할 것으로 판단됨 \par
1) 매출액 증가율 결정 방법: 대상 기업의 과거 매출액 평균 증가율(2017년∼2019년)과 동종 및 유사업종 매출액 평균 증가율의 산술 평균으로 결정함 \par
2) 동종 및 유사업종 매출액 증가율 \par
{\small\centering
\begin{tabular}{|c|c|c|c|c|}
\hline
\makecell{항목} & \makecell{단위} & \makecell{2017년} & \makecell{2018년} & \makecell{2019년} \\
\hline
\makecell{매출액 증가율} & \makecell{\%} & \makecell{4.92} & \makecell{4.82} & \makecell{5.24} \\
\hline
\end{tabular}\par}

2. 매출원가는 과거와 동일한 매출원가율을 적용함 \par
3. 판매비와 관리비는 향후에도 과거와 동일하게 매출액의 일정 비율만큼 발생할 것으로 봄 \par
4. 감가상각비는 2019년에는 110,000,000원이며 추정기간 동안 매년 5,000,000원씩 증가됨 \par
5. 향후 예상되는 자본적 지출액은 매출액의 3 \%가 소요될 것임 \par
6. 순운전자본 증감 \par
1) 대상 기업의 경우 추정 매출액 증감액에 운전자본 소요율을 곱하여 산출함 \par
(추정매출액$_{t}$ - 추정매출액$_{t-1}$) × 운전자본 소요율 \par
2) 운전자본 소요율은 한국은행 공시 재무정보를 이용한 회전율 등을 고려하여 대상 회사의 자료 등을 기준으로 산출하며, 추정기간 동안 동일하게 적용함 \par
\begin{center}
운전자본 소요율 $=\dfrac{1}{\text{매출채권회전율}}+\dfrac{1}{\text{재고자산회전율}}-\dfrac{1}{\text{매입채무회전율}}$
\end{center}
{\small\centering
\begin{tabular}{|c|c|c|c|}
\hline
\makecell{구분} & \makecell{매출채권회전율} & \makecell{재고자산회전율} & \makecell{매입채무회전율} \\
\hline
\makecell{회} & \makecell{8} & \makecell{10} & \makecell{20} \\
\hline
\end{tabular}\par}

7. 법인세 산정시 세율은 22 \%를 적용함 \par
< 자료 5 > 자기자본비용 관련 \par
1. 본 기업의 자본구조는 자기자본비율 40 \%, 타인자본비율 60 \%임 \par
2. 자기자본의 기회비용은 자본자산가격평가모델(CAPM법: Capital Asset Pricing Model)에 의함 \par
3. 무위험자산의 수익률(Rf)은 평균 5년 만기 국고채 수익률 등을 고려하여 3.5 \%, 시장의 기대수익률(E(Rm))은 12 \%로 가정함 \par
4. β계수는 최근 3년 유사업종에 속한 기업들의 β계수의 산술평균으로 함 \par
(식료품 제조업) \par
{\small\centering
\begin{tabular}{|c|c|}
\hline
\makecell{기준년도} & \makecell{기업베타(β)} \\
\hline
\makecell{2017년} & \makecell{0.9654} \\
\hline
\makecell{2018년} & \makecell{0.9885} \\
\hline
\makecell{2019년} & \makecell{0.9763} \\
\hline
\end{tabular}\par}

< 자료 6 > 타인자본비용 관련 \par
대상 기업의 재정상태 및 금융상환 가능성 등을 종합적으로 고려하여 대상 기업의 차입이자율을 7 \%로 결정함 \par
< 자료 7 > 특허권 평가 자료 \par
1. 특허권의 유효 잔존수명은 경제적 수명 잔존기간과 법적 잔존기간을 비교하여 결정하며, 산출된 유효 잔존수명은 연단위로 절사함 \par
2. 특허권의 경제적 수명 잔존기간은 아래의 자료로 산출함 \par
1) 경제적 수명기간 산출방법: 특허인용수명 × (1 + 영향요인 평점 합계 / 20) \par
2) 특허인용수명 \par
{\small\centering
\begin{tabular}{|c|c|c|c|c|}
\hline
\makecell{IPC} & \makecell{기술명} & \makecell{Q1} & \makecell{Q2(중앙값)} & \makecell{Q3} \\
\hline
\makecell{A23G} & \makecell{과자 등} & \makecell{5} & \makecell{9} & \makecell{13} \\
\hline
\end{tabular}\par}

※ 대상 특허의 특성 및 시장상황 등을 종합적으로 고려하여 대상 특허의 경제적 수명기간 산출에 적용할 특허인용수명은 중앙값으로 결정함 \par
3) 기술수명 영향요인 평가표 \par
{\small\centering
\begin{tabular}{|c|c|c|c|c|c|c|}
\hline
\multirow{2}{*}{\makecell{구분}} & \multirow{2}{*}{\makecell{세부요인}} & \multicolumn{5}{|c|}{\makecell{평점}} \\
\cline{3-7}
 &  & \makecell{-2} & \makecell{-1} & \makecell{0} & \makecell{1} & \makecell{2} \\
\hline
\multirow{5}{*}{\makecell{기술요인}} & \makecell{대체기술 출현가능성} &  &  &  & \makecell{v} &  \\
\cline{2-7}
 & \makecell{기술적 우월성} &  &  &  & \makecell{v} &  \\
\cline{2-7}
 & \makecell{유사ㆍ경쟁기술의 존재(수)} &  &  & \makecell{v} &  &  \\
\cline{2-7}
 & \makecell{모방 난이도} &  &  &  & \makecell{v} &  \\
\cline{2-7}
 & \makecell{권리 강도} &  &  & \makecell{v} &  &  \\
\hline
\multirow{5}{*}{\makecell{시장\newline{}요인}} & \makecell{시장 집중도 (주도기업 존재)} &  &  &  & \makecell{v} &  \\
\cline{2-7}
 & \makecell{시장경쟁의 변화} &  &  & \makecell{v} &  &  \\
\cline{2-7}
 & \makecell{시장경쟁강도} &  &  & \makecell{v} &  &  \\
\cline{2-7}
 & \makecell{예상 시장점유율} &  &  &  & \makecell{v} &  \\
\cline{2-7}
 & \makecell{신제품 출현빈도} &  &  &  & \makecell{v} &  \\
\hline
\end{tabular}\par}

3. 특허권은 물음 1에서의 “기업의 영업가치”에 해당 특허권의 기술기여도를 곱하는 방식으로 평가함 \par
< 자료 8 > 기술기여도 산출 관련 자료 \par
1. 결정방법: 기술기여도는 산업 특성을 반영하는 산업기술요소와 개별기술의 특성을 평가하는 개별기술강도의 곱으로 결정함 \par
2. 산업기술요소 \par
{\small\centering
\begin{tabular}{|c|c|c|c|c|}
\hline
\multicolumn{2}{|c|}{\makecell{표준산업분류코드}} & \makecell{최대무형자산가치\newline{}비율(\%)} & \makecell{기술자산비율\newline{}(\%)} & \makecell{산업기술요소\newline{}(\%)} \\
\hline
\makecell{C10} & \makecell{식료품 제조업} & \makecell{67.5} & \makecell{76.0} & \makecell{51.3} \\
\hline
\makecell{C28} & \makecell{전기장비제조업} & \makecell{90.4} & \makecell{75.3} & \makecell{68.1} \\
\hline
\end{tabular}\par}

※ 산업기술요소(\%) = 최대무형자산가치비율(\%) × 기술자산비율(\%) \par
3. 개별기술강도 \par
1) 기술성 \par
{\small\centering
\begin{tabular}{|c|c|c|c|c|c|c|}
\hline
\multirow{2}{*}{\makecell{구분}} & \multirow{2}{*}{\makecell{평가지표}} & \multicolumn{5}{|c|}{\makecell{평점}} \\
\cline{3-7}
 &  & \makecell{1} & \makecell{2} & \makecell{3} & \makecell{4} & \makecell{5} \\
\hline
\multirow{10}{*}{\makecell{기술성}} & \makecell{혁신성} &  &  &  & \makecell{v} &  \\
\cline{2-7}
 & \makecell{파급성} &  &  &  & \makecell{v} &  \\
\cline{2-7}
 & \makecell{활용성} &  &  & \makecell{v} &  &  \\
\cline{2-7}
 & \makecell{전망성} &  &  & \makecell{v} &  &  \\
\cline{2-7}
 & \makecell{차별성(독창성)} &  &  &  & \makecell{v} &  \\
\cline{2-7}
 & \makecell{대체성} &  &  &  & \makecell{v} &  \\
\cline{2-7}
 & \makecell{모방용이성} &  &  & \makecell{v} &  &  \\
\cline{2-7}
 & \makecell{진부화가능성(기술수명)} &  &  & \makecell{v} &  &  \\
\cline{2-7}
 & \makecell{권리범위} &  &  &  & \makecell{v} &  \\
\cline{2-7}
 & \makecell{권리 안정성} &  &  &  & \makecell{v} &  \\
\hline
\end{tabular}\par}

2) 사업성 \par
{\small\centering
\begin{tabular}{|c|c|c|c|c|c|c|}
\hline
\multirow{2}{*}{\makecell{구분}} & \multirow{2}{*}{\makecell{평가지표}} & \multicolumn{5}{|c|}{\makecell{평점}} \\
\cline{3-7}
 &  & \makecell{1} & \makecell{2} & \makecell{3} & \makecell{4} & \makecell{5} \\
\hline
\multirow{10}{*}{\makecell{사업성}} & \makecell{수요성} &  &  &  & \makecell{v} &  \\
\cline{2-7}
 & \makecell{시장진입성} &  &  &  & \makecell{v} &  \\
\cline{2-7}
 & \makecell{생산용이성} &  &  & \makecell{v} &  &  \\
\cline{2-7}
 & \makecell{시장점유율 영향} &  &  & \makecell{v} &  &  \\
\cline{2-7}
 & \makecell{경제적 수명} &  &  &  & \makecell{v} &  \\
\cline{2-7}
 & \makecell{매출 성장성} &  &  & \makecell{v} &  &  \\
\cline{2-7}
 & \makecell{파생적 매출} &  &  & \makecell{v} &  &  \\
\cline{2-7}
 & \makecell{상용화 요구시간} &  &  & \makecell{v} &  &  \\
\cline{2-7}
 & \makecell{상용화 소요자본} &  &  & \makecell{v} &  &  \\
\cline{2-7}
 & \makecell{영업 이익성} &  &  &  & \makecell{v} &  \\
\hline
\end{tabular}\par}

3) 개별기술강도 \par
개별기술강도(\%) = (기술성 점수 합산 + 사업성 점수 합산) / 100 \par
< 자료 9 > 영업권 평가 자료 \par
1. 영업권은 물음 1에서의 “기업의 영업가치(영업관련 기업가치)”에서 영업투하자본을 차감하는 방법으로 평가하되, 물음 2에서 평가된 특허권도 차감함 \par
2. 제시된 재무상태표를 기준으로 영업투하자본을 산출함 \par
< 자료 10 > 기타 \par
1. 기업가치는 “기업의 영업가치”와 비영업용자산으로 구성됨 \par
2. 연도별 매출액과 “기업의 영업가치”, 특허권 평가액, 영업권 평가액은 십만단위에서 반올림함 \par
3. 매출액 증가율을 제외한 모든 율은 백분율로 소수점 이하 셋째자리에서 반올림하여 백분율로 소수점 이하 둘째자리까지 표시함 \par
\end{longquestion}
